\chapter{Appendix}

\section{Source Code}
The complete source code of \textit{Micorrizas} is available in a GitHub repository. This repository includes the Unity project files, scripts, scenes and resources used during development, allowing the implementation of the prototype to be reviewed.

Repository: \url{https://github.com/Salix8/Micorrizas_JardinSentidos_TFG_Videojuegos.git}

\section{Mobile Device Simulation}
This appendix addresses the limitations found when attempting to simulate different mobile devices using a computer. During development, an attempt was made to test the geolocation system from a desktop environment, but this was not fully possible because computers do not usually provide the same sensors or location services as mobile devices.

On mobile platforms, there is a real location service available through the operating system and the device hardware. Unity provides access to this service through its location API. The LocationService.Start() function starts location updates, and the current position can then be accessed through Input.location.lastData. In Unity’s documentation, Input.location is described as a property used to access the location of handheld devices, that is, mobile devices.

In addition, on Android, when an application uses LocationService, Unity adds the ACCESS\_FINE\_LOCATION permission and requests location access from the user the first time the service is used. This behaviour is specific to mobile platforms.

On a PC, this process is different. A computer normally does not have GPS hardware or the same permission system used by Android or iOS. Although a PC may provide an approximate location based on IP address, Wi-Fi networks or operating system settings, this is not equivalent to reliable mobile geolocation. For this reason, the geolocation features of the project had to be validated using real mobile devices rather than relying only on computer-based simulation.

\section{Detailed technical architecture}
\label{sec:detailed-technical-architecture}
The project’s architecture has been designed in a modular way, separating multiplayer connectivity, co-op session management, geolocation acquisition, spatial representation on the map and mission-specific logic. This separation allows each subsystem to maintain a specific responsibility, with communication between them taking place via state structures.

The entry point for the co-operative system is the RelayConnectionService.cs connection service (line 15). This component is responsible for initialising Unity’s services, authenticating the user anonymously and configuring the network transport via Unity Relay. When a player creates a session, the service requests an allocation from Relay, obtains a join code and configures UnityTransport with the necessary data so that NetworkManager can start in host mode (RelayConnectionService.cs, line 121). When a player joins, the process is similar, but uses the received code to retrieve the remote allocation and start the client (RelayConnectionService.cs, line 172). In this way, Relay handles connectivity between devices, whilst Netcode for GameObjects takes over management of the network session and state replication.

Built on top of this infrastructure is the session’s main coordinator, `CoopSessionCoordinator.cs` (line 9), which acts as the authoritative hub of the co-operative flow. This component inherits from `NetworkBehaviour` and keeps various global session elements synchronised: the current game phase via a `NetworkVariable<CoopGamePhase>`, the active mission, and the list of registered players via a `NetworkList<ulong>` in `CoopSessionCoordinator.cs` (line 26). The coordinator is responsible for determining whether the group is in the Lobby, WorldMap, MiniGame or SessionSummary, as well as for orchestrating synchronised scene transitions for all clients via NetworkManager.SceneManager.LoadScene(...) in CoopSessionCoordinator.cs (line 387). Furthermore, it rebuilds player slots when connections are established or lost, so that other systems can determine which index each participant occupies within the session (CoopSessionCoordinator.cs, line 356).

The lobby interface does not contain any direct network logic, but acts as a presentation layer on top of these services. In MultiplayerMenuController.cs (line 7), the buttons to create a session, join by code, start a game or leave a session invoke operations from the RelayConnectionService and the CoopSessionCoordinator. This separation allows the UI to simply reflect the session’s state without duplicating business logic.

The geolocation subsystem begins in DeviceGpsService.cs (line 9), which encapsulates access to the device’s GPS via Input.location. This component requests permissions, checks whether the operating system’s location service is enabled, starts the sensor with configurable accuracy and distance, and manages timeouts and retries during initialisation (DeviceGpsService.cs, line 85). Each reading is converted into an immutable DeviceGpsReading object, which contains latitude, longitude, altitude, horizontal accuracy, service status, fix availability and timestamps (DeviceGpsService.cs, line 275). The service does not publish readings continuously and uncontrollably, but only when it detects a significant change via `HasMeaningfulChange(...)`, thereby reducing noise and unnecessary traffic (DeviceGpsService.cs, line 223).

The interaction between the local GPS and the cooperative system takes place via CoopGpsMarkerController.cs (line 10). This controller subscribes to the ReadingUpdated event of the DeviceGpsService, stores the latest local reading and immediately updates the player’s own marker on the map in CoopGpsMarkerController.cs (line 65). Furthermore, at regular intervals, it publishes this reading over the network via gpsStateSync.SubmitLocalReading(...) in CoopGpsMarkerController.cs (line 120). This architectural decision is important because it decouples the capture of the device’s position from its distribution in the multiplayer environment: DeviceGpsService merely measures, whilst CoopGpsMarkerController decides when that measurement is used for the co-operative environment.

Position synchronisation between players is handled in CoopGpsStateSync.cs (line 7). This component, which is also derived from NetworkBehaviour, maintains a NetworkList<CoopPlayerGpsState> containing the GPS status of all participants (CoopGpsStateSync.cs, line 9). When a client wishes to share its position, it calls SubmitLocalReading(...), which packages the data from DeviceGpsReading into a CoopPlayerGpsState and executes a ServerRpc to the host (CoopGpsStateSync.cs, line 49). On the server, SubmitGpsStateServerRpc(...) identifies the sender by SenderClientId and updates or inserts their entry into the shared list in CoopGpsStateSync.cs (line 94). From that point onwards, NGO automatically replicates the list to the other clients. Therefore, the actual flow is: local GPS → structured reading → periodic publication → ServerRpc → NetworkList update → replication to the other players.

These coordinates are represented spatially using ArcGIS. The ArcGISMapCoordinateProjector.cs component (line 6) receives an ArcGISLocationComponent and assigns it an ArcGISPoint in the WGS84 system using longitude, latitude and altitude (ArcGISMapCoordinateProjector.cs, line 16). CoopGpsMarkerController uses this projector in UpdateMarker(...) to place a visible marker for each player, whether local or remote, within the 3D map (CoopGpsMarkerController.cs, line 178). In this way, the synchronised GPS status is converted into a spatial position that can be navigated and utilised by the rest of the game’s systems.

The same marker controller also links geolocation to the players’ visual profiles. To do this, it relies on `CoopPlayerProfileSync.cs` (line 9), which synchronises each client’s `avatarId` and display name over the network using another `NetworkList` in `CoopPlayerProfileSync.cs` (line 15). The local profile is sent to the server using SubmitLocalProfile() and SubmitProfileServerRpc(...) in CoopPlayerProfileSync.cs (line 81). Subsequently, `CoopGpsMarkerController` queries these profiles to determine which appearance each marker should use via `ResolveDesiredAvatarId(...)` in `CoopGpsMarkerController.cs` (line 546). In this way, a player’s remote position is shared not only as coordinates, but also as a visual identity within the co-operative map.

The relationship between geolocation and access to missions is reflected in the map’s zone system. Each zone is described by `CoopMinigameZoneDefinition.cs` (line 5), which defines which mission it represents, its display name, the maximum acceptable GPS accuracy threshold, and the spatial tolerance considered valid for being ‘within’ the zone. CoopMinigameZoneTriggerController.cs (line 11) works with these definitions, periodically checking whether all registered players are simultaneously within the same valid zone (CoopMinigameZoneTriggerController.cs, line 99). To do this, it does not use latitude and longitude directly, but instead queries CoopGpsStateSync for each player’s latest status and CoopGpsMarkerController for their position, already projected into the Unity world, using TryGetMarkerWorldPosition(...) in CoopMinigameZoneTriggerController.cs (line 162). If all participants meet the spatial condition and maintain that position during a countdown, the controller requests the CoopSessionCoordinator to load the corresponding minigame using StartMiniGame(...) in CoopMinigameZoneTriggerController.cs (line 158). This is one of the most important relationships in the architecture: the GPS feeds the co-op system and, via the map, acts as a condition for accessing the playable content.

The global session progress layer is located in CoopSessionProgressSync.cs (line 9). This component maintains a NetworkList<CoopMinigameProgressNetworkState> that records which missions have been completed, in what order, and with what average score (CoopSessionProgressSync.cs, line 15). The session coordinator uses this to prevent a minigame that has already been completed from being launched again within the same game, and also to decide whether, after a mission, the group should return to the map or proceed to the final summary scene (CoopSessionCoordinator.cs, line 216). The result of each mission is recorded using `TryRegisterResultServer(...)`, ensuring that the overall progress state remains synchronised across all clients (CoopSessionProgressSync.cs, line 93).

The missions themselves follow a common architecture based on inheritance and reusability. The abstract class CooperativeMinigameBase.cs (line 9) defines the standard lifecycle of all missions: tutorial, waiting for players, active phase and results. To do this, it uses several NetworkVariables, including the phase status, the published result and the list of players who have already completed the tutorial CooperativeMinigameBase.cs (line 14). Each specific mission implements `InitializeMinigameServer()` to set up its internal state and can call `PublishResultServer(...)` to communicate the result to both the local minigame flow and the global session progress (CooperativeMinigameBase.cs, line 208). Thanks to this common foundation, all minigames share the same behaviour regarding the tutorial, blocking, progression following results and return to the map, without having to reimplement that flow.

A clear example of this pattern is AudioWordConsensusMinigameSession.cs (line 11). This class inherits from CooperativeMinigameBase and adds its own NetworkVariables to control the active round, sound emitter, disabled words, time remaining, errors made and shared message (AudioWordConsensusMinigameSession.cs, line 17). The server initialises the game, generates the round plan, processes client responses via RPCs and, when the mission ends, publishes a `MinigameResultData` which is returned to the session coordinator and the global progress system (AudioWordConsensusMinigameSession.cs, line 310; AudioWordConsensusMinigameSession.cs, line 560). The rest of the project’s missions follow a similar structure: their own configuration, synchronised session, network models, auxiliary services and specific UI.

In addition to the gameplay logic, there are several supporting components that reinforce the architecture. CoopInformationPresenter.cs (line 5) uses the coordinator’s slot mapping to display different variants of information depending on the local player. MapSceneEntryGuard.cs (line 7) prevents invalid direct entries to the map outside the co-op flow. CoopMinigameLauncherUIController.cs (line 11) dynamically builds the panel of available minigames based on settings, the current phase and shared progress. Finally, GpsDebugPanelUIController.cs (line 8) provides a technical debugging view that allows real-time verification of the local reading, the latest synchronised GPS status, the ArcGIS projection and the final position of the marker in the world, which is useful for validating the entire chain from the sensor to the spatial representation.

The project’s architecture can be summarised as a layered data flow. First, the device obtains a local GPS reading via DeviceGpsService. Next, CoopGpsMarkerController takes that reading and sends it periodically to CoopGpsStateSync, which replicates it to the other clients via NGO. Next, CoopGpsMarkerController and ArcGISMapCoordinateProjector transform this synchronised state into visible markers on the map. These markers then serve as the basis for the spatial logic of CoopMinigameZoneTriggerController, which detects whether all players are within the same zone and, if so, requests the CoopSessionCoordinator to start the corresponding mission. Once inside the mission, the specific session inherits the common workflow from `CooperativeMinigameBase`, publishes its result and returns it to the coordinator via `CoopSessionProgressSync`, thereby closing the loop between the session, geolocation, map and cooperative progress.